
\vspace{-0.8cm}

This chapter describes the methodology and implementation process of the \textbf{Automated Car Parking System}. It explains the development steps, tools, data handling, algorithm, flowchart, and main implementation modules.

\section{Methodological Framework}

The project was developed using a step-by-step embedded system approach. First, the parking problem and required features were identified. Then suitable components were selected based on cost, availability, and compatibility with Arduino UNO. After that, the circuit and robotic arm gate structure were designed. Each module was tested separately and finally integrated into one complete system.

The methodology followed in this project is given below:

\begin{enumerate}
    \item Identify the parking management problem and define objectives.
    \item Select required hardware components.
    \item Design the circuit and two-link robotic arm gate.
    \item Develop Arduino code for sensors, RFID, gate control, LCD, and Wi-Fi communication.
    \item Test each module individually.
    \item Integrate all modules into the final prototype.
    \item Test the system under entry and exit conditions.
\end{enumerate}

\section{Tools, Libraries, and Technologies Used}

The system was developed using both hardware and software tools. Arduino UNO was used as the main controller, and Arduino IDE was used for writing and uploading the program.

\begin{table}[H]
\centering
\small
\caption{Tools, Libraries, and Technologies Used}
\label{tab:tools_technologies}
\begin{tabular}{p{4cm} p{8.5cm}}
\toprule
\textbf{Tool / Technology} & \textbf{Purpose} \\
\midrule
Arduino IDE & Writing, compiling, and uploading Arduino code. \\
Arduino UNO & Main microcontroller of the system. \\
C/C++ & Programming language for Arduino. \\
IR Sensors & Vehicle detection at entry and exit points. \\
MFRC522 RFID Module & RFID card-based authentication. \\
Servo Motor & Controls the two-link robotic arm gate. \\
16x2 I2C LCD & Displays slot availability and system messages. \\
ESP-01S Wi-Fi Module & Sends parking data to the online server. \\
Firebase / Online Database & Stores vehicle entry and exit records. \\
\bottomrule
\end{tabular}
\normalsize
\end{table}

The main Arduino libraries used are \texttt{Wire.h}, \texttt{LiquidCrystal\_I2C.h}, \texttt{Servo.h}, \texttt{SPI.h}, and \texttt{MFRC522.h}. These libraries support I2C communication, LCD control, servo movement, SPI communication, and RFID reading.

\section{Data Collection and Processing}

This project does not use a machine learning dataset. It uses real-time data from hardware modules. The main data includes IR sensor status, RFID UID, slot count, event type, timestamp, and owner information.

\begin{table}[H]
\centering
\small
\caption{System Data Description}
\label{tab:data_description}
\begin{tabular}{p{4cm} p{8.5cm}}
\toprule
\textbf{Data Type} & \textbf{Description} \\
\midrule
IR Sensor Data & Detects vehicle presence at entry or exit. \\
RFID UID & Stores the unique ID of the scanned RFID card. \\
Slot Count & Maintains available parking slots. \\
Event Type & Identifies entry or exit operation. \\
Timestamp & Stores the time of the parking event. \\
Owner Information & Identifies the user linked with the RFID card. \\
\bottomrule
\end{tabular}
\normalsize
\end{table}

The Arduino processes the collected data. RFID UID is converted into a readable string, trimmed, and matched with authorized UID values. The system also prepares a timestamp and checks the slot count to avoid invalid updates. For every valid entry or exit, a JSON packet is created and sent online through the ESP-01S Wi-Fi module.

\section{Model Development / Algorithm Design}

The project does not use a machine learning model. It follows rule-based control logic. The Arduino continuously checks the entry and exit IR sensors. If a vehicle is detected at the entry point, the system checks slot availability and requests RFID authentication. If the RFID card is valid, the gate opens and the slot count is updated.

\subsection{System Flowchart}

Fig.~\ref{fig:flowchart_diagram} shows the working flow of the Automated Car Parking System. It includes hardware initialization, vehicle detection, slot checking, RFID verification, robotic arm gate control, LCD update, and online data transmission.

\begin{figure}[H]
    \centering
    \includegraphics[height=0.72\textheight]{flowchart_of_automated_car_parking_system.png}
    \caption{Flowchart of Automated Car Parking System}
    \label{fig:flowchart_diagram}
\end{figure}

The simplified algorithm is given below:

\begin{algorithm}[H]
\caption{System Operation Algorithm}
\begin{algorithmic}[1]
\State Initialize LCD, RFID, servo motor, IR sensors, and Wi-Fi module
\State Set maximum parking slot value
\While{system is active}
    \If{entry sensor detects a car}
        \If{slot is available}
            \State Request RFID card scan
            \If{RFID card is valid}
                \State Open robotic arm gate
                \State Decrease slot count
                \State Send entry data to online server
                \State Close robotic arm gate
            \Else
                \State Display access denied message
            \EndIf
        \Else
            \State Display parking full message
        \EndIf
    \EndIf

    \If{exit sensor detects a car}
        \State Open robotic arm gate
        \State Increase slot count
        \State Send exit data to online server
        \State Close robotic arm gate
    \EndIf

    \State Update LCD display
\EndWhile
\end{algorithmic}
\end{algorithm}

\section{System Implementation / Modules Description}

The system was implemented by dividing it into functional modules. Each module performs a specific task and communicates with the Arduino UNO.

\subsection{Vehicle Detection Module}

Two IR sensors are used to detect vehicles at the entry and exit points. When a vehicle is detected, the sensor sends a digital signal to the Arduino and starts the entry or exit process.

\subsection{RFID Authentication Module}

The RFID module checks whether the user is authorized. If the scanned card matches a stored UID, the system allows entry. Otherwise, the LCD displays a wrong card or access denied message.

\subsection{Robotic Arm Gate Module}

The gate is controlled by a servo motor connected to a two-link robotic arm. When entry or exit is allowed, the servo moves the arm to the open position. After the vehicle passes, the arm returns to the closed position. Smooth movement is created by changing the servo angle step by step.

\subsection{LCD Display Module}

The 16x2 I2C LCD displays welcome messages, RFID scanning messages, wrong card alerts, parking full messages, gate status, and available slot count. This helps users understand the current system condition.

\subsection{IoT Communication Module}

The ESP-01S Wi-Fi module sends parking data to the online server using AT commands. For each valid event, the Arduino creates a JSON message and sends it through an HTTP POST request.

\begin{lstlisting}
{
  "uid": "F6 A9 23 03",
  "event": "entry",
  "owner": "MRJ",
  "timestamp": "00:05:23"
}
\end{lstlisting}

\section{Hardware and Software Requirements}

\begin{table}[H]
\centering
\small
\caption{Hardware and Software Requirements}
\label{tab:hardware_software_req}
\begin{tabular}{p{4cm} p{8.5cm}}
\toprule
\textbf{Category} & \textbf{Requirements} \\
\midrule
Controller & Arduino UNO \\
Sensors & Two IR sensors for entry and exit detection \\
Authentication & MFRC522 RFID reader and RFID cards \\
Actuator & Servo motor with two-link robotic arm mechanism \\
Display & 16x2 I2C LCD display \\
Communication & ESP-01S Wi-Fi module \\
Power Supply & 5V supply for Arduino and connected modules \\
Programming Tool & Arduino IDE \\
Programming Language & C/C++ for Arduino \\
Libraries & Wire, LiquidCrystal\_I2C, Servo, SPI, MFRC522 \\
Online Platform & Online server or database for parking data storage \\
\bottomrule
\end{tabular}
\normalsize
\end{table}

\section{Summary}

This chapter explained the methodology and implementation of the Automated Car Parking System. It described the development steps, tools, system data, processing method, flowchart, algorithm, implementation modules, and hardware-software requirements. The final implementation combines vehicle detection, RFID authentication, robotic arm gate control, LCD display, and IoT-based data transmission.

